\documentclass[oneside]{estilo}

% Importações de pacotes
\usepackage[alf, abnt-emphasize=bf, recuo=0cm, abnt-etal-cite=2, abnt-etal-list=0]{abntex2cite}  % Citações padrão ABNT
\usepackage[utf8]{inputenc}                         % Acentuação direta

\usepackage[T1]{fontenc}                            % Codificação da fonte em 8 bits
\usepackage{graphicx}                               % Inserir figuras
\usepackage{amsfonts, amssymb, amsmath}             % Fonte e símbolos matemáticos
\usepackage{booktabs}                               % Comandos para tabelas
\usepackage{verbatim}                               % Texto é interpretado como escrito no documento
\usepackage{multirow, array}                        % Múltiplas linhas e colunas em tabelas
\usepackage{indentfirst}                            % Endenta o primeiro parágrafo de cada seção.
\usepackage{microtype}                              % Para melhorias de justificação?
\usepackage[algoruled, portuguese]{algorithm2e}     % Escrever algoritmos
\usepackage{float}                                  % Utilizado para criação de floats
%\usepackage{times}                                 % Usa a fonte Times


\usepackage{color}
\usepackage{chronosys}
\usepackage{pdflscape}
\usepackage{tabu}
\usepackage{makecell}
\usepackage{nameref}
\usepackage{listings}
\usepackage{enumitem}
\usepackage{chngcntr}
\usepackage{colortbl}
\usepackage[htt]{hyphenat}
\usepackage{mdframed}

\mdfsetup{skipabove=1.5\topskip,skipbelow=0.5\topskip}
\mdfdefinestyle{pp}{%
	linecolor=black,
	outerlinewidth=.5pt,
	roundcorner=5pt,
	innertopmargin=8pt,
	innerbottommargin=0.5\baselineskip,
	usetwoside=false
	%backgroundcolor=gray!50!white}
}

%\usepackage{minitoc}
%\dominitoc% Initialization
\newcommand{\enumerateletter}[2]{\textbf{#1\textsubscript{#2}}}

% Inclui o preâmbulo do documento
%
% Documento: Preâmbulo
%
\usepackage{lmodern}
\usepackage [latin1]{inputenc}
\renewcommand{\familydefault}{\rmdefault}
\renewcommand{\lstlistingname}{Código}

\counterwithin{quadro}{chapter}
\counterwithin{figure}{chapter}
\counterwithin{table}{chapter}
% Altera o rótulo a ser usando no elemento pré-textual "Lista de código"
\renewcommand{\lstlistlistingname}{Lista de códigos}

% Configura a ``Lista de Códigos'' conforme as regras da ABNT (para abnTeX2)
\begingroup\makeatletter
\let\newcounter\@gobble\let\setcounter\@gobbletwo
\globaldefs\@ne \let\c@loldepth\@ne
\newlistof{listings}{lol}{\lstlistlistingname}
\newlistentry{lstlisting}{lol}{0}
\endgroup

\renewcommand{\cftlstlistingaftersnum}{\hfill--\hfill}

\let\oldlstlistoflistings\lstlistoflistings
\renewcommand{\lstlistoflistings}{%
	\begingroup%
	\let\oldnumberline\numberline%
	\renewcommand{\numberline}{\lstlistingname\space\oldnumberline}%
	\oldlstlistoflistings%
	\endgroup}

\lstdefinestyle{master-thesis}{
basicstyle=\ttfamily,
numbers=left,
numberstyle=\tiny,
frame=tb,
columns=fullflexible,
showstringspaces=false
}
\lstset{style=master-thesis}

\instituicao{Universidade Federal de Minas Gerais}
\localapresentacao{Belo Horizonte - Minas Gerais}
\programa{Graduação em Engenharia de Sistemas}
\modalidade{}
\nomeautor{André Versiani de Mattos}
\titulotb{Caracterização e Extração de Atributos de Faltas de Alta Impedância em Vegetação: Aplicação em Dados Reais}
%\subtitulo{Subtítulo do trabalho}
\anoapresentacao{2026}
\grau{Bacharel em Engenharia de Sistemas}
\dataapresentacao{2026}
\mesapresentacao{Dezembro}

%Dados Orientador
\orientador{Gabriela Nunes}
\instOrientador{Cesar School}
\titulacaoorientador{Prof.ª Dra.}

%Dados Coorientador
% \coorientador{Professor Fulano}
% \instCoorientador{Instituição de Ensino}
% \titulacaocoorientador{Prof. Dr.}

\usepackage[a4paper, top=3.0cm, bottom=2.0cm, left=3.0cm, right=2.0cm]{geometry} %teste margem
%\usepackage{titlesec} %teste margem

\newcommand{\blue}[1]{\textcolor{blue}{#1}} % proposed
\newcommand{\red}[1]{\textcolor{red}{#1}} % to be removed
\newcommand{\green}[1]{\textcolor{green}{#1}} % to be discussed?
%\newcommand{\yellow}[1]{\textcolor{yellow}{#1}} % ????
\newcommand{\magenta}[1]{\textcolor{magenta}{#1}} % warning - requires attention

\settocdepth{subsection}

% Define as cores dos links e informações do PDF
\makeatletter
\hypersetup{
    portuguese,
    colorlinks,
    linkcolor=blue,
    citecolor=blue,
    filecolor=black,
    urlcolor=black,
    breaklinks=true,
    pdftitle={\titulotb{}},
    pdfauthor={\nomeautor{}},
    pdfsubject={\imprimirpreambulo},
    pdfkeywords={metamorphic testing, oracle problem, software testing}
}
\makeatother

% Redefinição de labels
\renewcommand{\algorithmautorefname}{Algoritmo}
\def\equationautorefname~#1\null{Equa\c c\~ao~(#1)\null}

% Cria o índice remissivo
\makeindex

% Início do documento
\begin{document}

    % Retira espaço extra obsoleto entre as frases.
    \frenchspacing
    % Elementos pré textuais
    \pretextual
    \include{./01-elementos-pre-textuais/capa}              % Capa
    %
% Documento: FOLHA DE ROSTO
%

\makeatletter
\begin{folhaderosto}
	\thispagestyle{empty}%limpa estilo da pagina
	
    \begin{center}
    
		\small\textbf{\expandafter\uppercase\expandafter{\imprimirnomeautor}}\\
		\vspace*{5.2 cm}%Espaço entre linhas
		\normalsize\textbf{\expandafter\uppercase\expandafter{\imprimirtitulotb}}\\
		
    \end{center}
	
	\vspace*{2.35 cm}%Espaçamento entre linhas
		    \large%tamanho da fonte 
    		\hfill%Estica horizontamente  com espaços
	    	\begin{minipage}{8 cm}%Minipagina
	    		\begin{small} %Muda tamanho da fonte
	    		\setlength{\baselineskip}{0.8\baselineskip}
				
				
		    	{Dissertação apresentada junto ao programa de {\textbf{\imprimirprograma}}
		    	do {\textbf{\imprimirinstituicao}},
		    	como requisito parcial à obtenção do título de
		    	{\textbf{\imprimirgrau}}.}\\{
		    	}\\\textbf{Orientadora}:\\ \\
		    	{\imprimirtitulacaoorientador }{ }{\imprimirorientador.}\\{
		    	}
				
				
				\end{small} %Muda tamanho da fonte
		    \end{minipage}%%Minipagina
		    	
		    \vspace*{6.5 cm}%Espaçamento entre linhas
		    
		    \begin{center} %Alinhamento centralizado
		    	\normalsize %Muda tamanho da fonte
	    		\textbf{\imprimirlocalapresentacao, \imprimirdataapresentacao}
	    	\end{center}%Alinhamento centralizado

\end{folhaderosto}
\makeatother
        % Folha de rosto
    % %
% Documento: Agradecimento
%

\begin{agradecimento}	
		\noindent Texto agradecendo pessoas que lhe ajudaram durante a produção do seu trabalho, seja com apoio técnico ou emocional.
\end{agradecimento}

    % Agradecimentos
    % \include{./01-elementos-pre-textuais/resumoPt}          % Resumo na língua vernácula
    % \include{./01-elementos-pre-textuais/resumoEn}          % Resumo em língua estrangeira
    % \include{./01-elementos-pre-textuais/listaFiguras}      % Lista de figuras
    % \include{./01-elementos-pre-textuais/listaTabelas}      % Lista de tabelas
    % \include{./01-elementos-pre-textuais/listaQuadros}      % Lista de quadros
    % \include{./01-elementos-pre-textuais/listaSiglas}       % Lista de abreviaturas e siglas
    % \include{./01-elementos-pre-textuais/listaSimbolos}     % Lista de símbolos
    \sumario
\begin{OnehalfSpacing}

    % Elementos textuais
    \textual
        %
% Documento: Introdução
%
\chapter{Introdução}\label{chap:introducao}

A energia elétrica constitui um insumo essencial para o funcionamento da sociedade contemporânea, sustentando atividades produtivas, serviços públicos e aplicações críticas em áreas como saúde, comunicação e infraestrutura. Nos sistemas de distribuição, a continuidade do fornecimento e a confiabilidade da operação dependem da capacidade de identificar perturbações que, embora nem sempre apresentem correntes elevadas, podem comprometer a segurança e a integridade da rede.

Nesse contexto, as falhas de alta impedância (FAIs) são reconhecidas, desde os trabalhos clássicos da década de 1980, como um dos problemas mais persistentes e difíceis enfrentados pela proteção de sistemas elétricos. \cite{1,2} Essas ocorrências foram descritas como situações em que um condutor energizado entra em contato com meios de alta resistência, como solo seco, asfalto ou outros materiais semelhantes, produzindo correntes insuficientes para a atuação confiável das proteções convencionais por sobrecorrente. \cite{3} Além disso, a literatura demonstra que as FAIs apresentam características complexas, com comportamento não linear, assimetria e assinaturas dependentes das condições de contato e do ambiente, o que dificulta sua detecção e localização. \cite{4,5}

Com o avanço das pesquisas a partir dos anos 2000, diferentes métodos passaram a ser propostos para melhorar a identificação dessas ocorrências, incluindo abordagens baseadas em componentes de alta frequência, harmônicos, inter-harmônicos e análise no domínio do tempo e da frequência. \cite{5,6,7,8} Esses estudos indicam que, embora haja progresso na modelagem e na detecção, o problema permanece relevante para a engenharia de proteção, sobretudo quando se consideram as variações reais do sistema e as dificuldades de implementação em campo. \cite{6,7}

Quando associadas à vegetação, as FAIs assumem importância ainda maior, pois podem prolongar a ocorrência de arco elétrico e aumentar o risco de ignição de materiais combustíveis. A literatura recente destaca que as falhas relacionadas à vegetação podem causar danos à rede, dificuldades operacionais e até incêndios, o que as torna particularmente críticas em áreas com arborização próxima às linhas aéreas. \cite{8,9}

Os dados australianos utilizados neste trabalho \cite{vic_bushfire_report} foram obtidos a partir do programa experimental Vegetation Conduction Ignition Test Program, desenvolvido no âmbito do Powerline Bushfire Safety Program do governo de Victoria, após os incêndios florestais de 2009. Esse programa foi criado para aprofundar o entendimento dos processos de ignição em falhas envolvendo vegetação e para estabelecer uma base de referência de assinaturas elétricas de falha, com o objetivo de apoiar o desenvolvimento de métodos mais eficientes de detecção. O conjunto de dados é especialmente relevante porque contém registros reais de testes controlados, com múltiplas espécies vegetais, diferentes condições experimentais e sinais elétricos capturados em alta resolução, o que o torna adequado para a formulação de um problema de classificação.

Diante do exposto, o presente estudo propõe utilizar esses registros para construir um problema de classificação supervisionada, no qual as medições elétricas associadas às ocorrências de falha serão analisadas a fim de distinguir classes relacionadas à presença ou ausência de contato vegetação–linha, ou ainda categorias derivadas do tipo de vegetação, severidade da ignição ou estágio do evento.            % Introdução
        %
% Documento: Revisão Bibliografica
%
\chapter{Revisão da Literatura}\label{chap:revbibliografica}

\section{Faltas de Alta Impedância}

Definição, características elétricas (baixa corrente, arco, não linearidade), tipos de superfície de contato, e o caso especifico da vegetação.

\section{Métodos de Detecção de Falhas}

Revisão cronológica dos métodos clássicos (harmônicos, componentes de alta frequência) ate os mais recentes (wavelets, inter-harmônicos).

\section{Técnicas de Processamento de Sinais}

Transformada de Fourier, transformada wavelet...

\section{Algoritmos de Inteligência Artificial}

Árvores de decisão, SVM, redes neurais, XGBoost, etc...

\section{Base de Dados Australiana}

O Vegetation Conduction Ignition Test Program, o dataset VeHIF
    % \include{./02-elementos-textuais/disposicoes}           % Disposições
    % \include{./02-elementos-textuais/estrutura}  			% Estrutura    
    % \include{./02-elementos-textuais/ambientes}  			% Estrutura    
    % \include{./02-elementos-textuais/ilustracao}           	% Ilustrações
    % \include{./02-elementos-textuais/citacoes}            	% Citações
    % \include{./02-elementos-textuais/tecnicaref}            % Tecnica de referencias
    
\end{OnehalfSpacing}
    % Elementos pós textuais
    \postextual
    \include{./03-elementos-pos-textuais/referencias}       % Referências
    \begin{OnehalfSpacing}
%    \include{./03-elementos-pos-textuais/apendices}         % Apêndices
%    \include{./03-elementos-pos-textuais/anexos}            % Anexos
    \end{OnehalfSpacing}
    %\printindex                                             % Índice remissivo

\end{document}
